\documentclass[11pt,a4paper]{article}

% --- Core Typography & Geometry Packages ---
\usepackage[utf8]{inputenc}
\usepackage[T1]{fontenc}
\usepackage{lmodern}
\usepackage{geometry}
\geometry{top=1in, bottom=1in, left=0.85in, right=0.85in}

\usepackage{amsmath, amssymb}
\usepackage{booktabs}
\usepackage{tabularx}
\usepackage{enumitem}
\usepackage{graphicx}
\usepackage[space]{grffile}
\usepackage{float}
\usepackage{caption}
\usepackage{microtype}
\usepackage{listings}
\usepackage{titlesec}
\usepackage{setspace}
\setstretch{1.15}

% --- Strict URL Wrapping (Prevents Margin Overflow) ---
\usepackage{xurl}

% --- Pure Black Styling (No Colors - Professional Academic Format) ---
\usepackage[hidelinks]{hyperref}
\hypersetup{
    colorlinks=false,
    pdfborder={0 0 0}
}
\urlstyle{same}

% --- Code Listings Configuration (Monochrome) ---
\lstset{
    basicstyle=\ttfamily\footnotesize,
    breaklines=true,
    frame=single,
    rulecolor=\color{black},
    numbers=left,
    numberstyle=\tiny\color{black},
    keywordstyle=\bfseries,
    commentstyle=\itshape,
    stringstyle=\ttfamily,
    showstringspaces=false,
    tabsize=2,
    captionpos=b
}

% --- Section Heading Formatting ---
\titleformat{\section}{\large\bfseries}{\thesection.}{0.5em}{}[\titlerule]
\titleformat{\subsection}{\normalsize\bfseries}{\thesubsection}{0.5em}{}
\titleformat{\subsubsection}{\small\bfseries\itshape}{\thesubsubsection}{0.5em}{}

\begin{document}

% =========================================================================
% TITLE & AUTHORS (CENTERED AT TOP)
% =========================================================================
\begin{center}
    {\LARGE \textbf{GitLock: Trustless Code Escrow and Verifiable Software Delivery Protocol on BNB Smart Chain}}\\[1.2em]
    
    {\large \textbf{Academic Case Study Report}}\\[0.4em]
    {\normalsize Department of Computer Science and Engineering}\\[1.2em]
    
    \begin{tabular}{cc}
        \textbf{Edavalapati Aashiq} & \textbf{Pucha Tirupathi Reddy} \\
        Roll No: CB.SC.U4CSE23560 & Roll No: CB.SC.U4CSE23568 \\
    \end{tabular}\\[1.5em]
    
    {\small \textbf{Project Artifacts \& Live Deployment Links:}}\\[0.5em]
\end{center}

\noindent
\begin{tabularx}{\textwidth}{@{} >{\bfseries}p{4.2cm} X @{}}
    Frontend Application (Vercel): & \url{https://git-lock-delta.vercel.app/} \\
    \addlinespace[2pt]
    Backend Service (Cloud Run): & \url{https://gitlock-backend-197473198672.asia-southeast1.run.app/api} \\
    \addlinespace[2pt]
    Smart Contract Address: & \texttt{0x30CF8a5fa441D39e0604781C24ADeeCBcbfecD49} \\
    \addlinespace[2pt]
    Blockchain Explorer (BscScan): & \url{https://testnet.bscscan.com/address/0x30CF8a5fa441D39e0604781C24ADeeCBcbfecD49} \\
\end{tabularx}

\vspace{0.8em}
\hrule
\vspace{1em}

% =========================================================================
% ABSTRACT
% =========================================================================
\begin{abstract}
Traditional freelance software engineering and remote contractual development suffer from an inherent counterparty dilemma: clients fear paying upfront for incomplete or non-functional deliverables, while developers risk non-payment after relinquishing proprietary source code. Existing Web2 intermediaries introduce exorbitant platform fees (10\%--20\%), opaque arbitration processes, and custodial single-points-of-failure. 

This paper presents \textbf{GitLock}, a decentralized escrow protocol and verifiable delivery vault designed to eliminate counterparty risk without trusted human intermediaries. GitLock bridges off-chain software repository lifecycle management with on-chain cryptographic settlement on the BNB Smart Chain (BSC) Testnet. The system enforces an atomic, verifiable seven-step delivery protocol: source code archives are cryptographically hashed and hybrid-encrypted using AES-256-CBC, live deployments are verified against specific Git commit hashes, and contractual funds are locked within an immutable Solidity smart contract (\texttt{GitLockEscrow}). Payment release is strictly gated by an Elliptic Curve Digital Signature Algorithm (ECDSA) attestation signed by a trusted GitLock Oracle only after cryptographically verifying complete repository ownership transfer to the client via the GitHub REST API. This architecture guarantees non-custodial atomicity, zero escrow commissions, cryptographic proof of delivery, and automated dispute timeouts.
\end{abstract}

\vspace{0.5em}
\noindent\textbf{Keywords:} Smart Contracts, Blockchain Escrow, Verifiable Software Delivery, GitHub API, ECDSA Oracle Attestation, BNB Smart Chain, Hybrid Encryption.

\vspace{1em}

% =========================================================================
% 1. INTRODUCTION & BACKGROUND
% =========================================================================
\section{Introduction and Background}
The exponential expansion of decentralized software engineering and global remote freelancing has highlighted the limitations of conventional contract fulfillment mechanisms. In peer-to-peer digital labor, deliverables consist entirely of intangible intellectual property: source code, deployment scripts, and repository administration rights. Once source code is exposed or ownership is transferred, revocation is technically impossible.

Existing commercial freelance platforms attempt to resolve this dilemma using centralized custodial escrow accounts. While functional, centralized intermediaries exhibit critical vulnerabilities:
\begin{enumerate}[noitemsep]
    \item \textbf{Prohibitive Intermediation Rent:} Centralized platforms levy significant platform commissions (ranging from 10\% to 20\%) on transacted sums.
    \item \textbf{Subjective Arbitration:} When disputes occur, non-technical human adjudicators lack the programmatic tools required to verify whether commit trees, build artifacts, or deployment manifests satisfy specifications.
    \item \textbf{Payment Latency and Custodial Risk:} Settling fiat payments often requires 7 to 14 business days, exposing funds to bank freezes, regulatory holds, and administrative clawbacks.
\end{enumerate}

Blockchain technology enables programmable conditional value transfers via smart contracts. By encoding contract settlement terms into deterministic bytecode executed on a distributed virtual machine, financial execution becomes strictly bound to verifiable state transitions. GitLock realizes this paradigm for software engineering by coupling automated GitHub repository inspection with an on-chain, non-custodial escrow contract.

% =========================================================================
% 2. PROBLEM STATEMENT
% =========================================================================
\section{Problem Statement}
In distributed freelance software delivery, clients and independent developers encounter the classic game-theoretic \textit{Hold-up Problem}:
\begin{itemize}[noitemsep]
    \item \textbf{The Developer's Dilemma:} If the developer transfers full administrative repository ownership or delivers decrypted production code prior to receiving payment, an unscrupulous client can terminate communication, dispute payments, or deploy the code without compensation.
    \item \textbf{The Client's Dilemma:} If the client deposits funds directly to the developer's wallet before receiving administrative ownership of the GitHub repository, the developer may deliver corrupted files, broken branches, or disappear entirely.
    \item \textbf{Verification Gap:} In modern web development, source code compilation alone does not prove delivery. A deliverable must be verified in a functional staging environment and tied to an immutable commit hash.
\end{itemize}

The technical challenge is to design an automated protocol that achieves \textbf{conditional atomicity}: the developer is guaranteed financial compensation if and only if repository ownership is verified as transferred to the client, while the client is guaranteed complete receipt of the verified codebase and repository rights or an automatic refund upon developer timeout.

% =========================================================================
% 3. PROPOSED SOLUTION & ARCHITECTURE
% =========================================================================
\section{Proposed Solution and System Architecture}
GitLock resolves this dilemma by decoupling code verification, staging demonstration, repository administration transfer, and on-chain monetary settlement into four coordinated tiers:
\begin{enumerate}[noitemsep]
    \item \textbf{Client Presentation Layer (Next.js 14):} Non-custodial Web3 interface enabling wallet connectivity, interactive project negotiation, code locking, and deal tracking.
    \item \textbf{Cryptographic Oracle \& API Backend (Express.js \& Node.js):} Off-chain verifier that performs cryptographic commit checksum verification, tracks Vercel staging deployment statuses, inspects GitHub repository ownership transitions, and issues ECDSA attestation signatures.
    \item \textbf{Relational Metadata \& Persistence Layer (Supabase PostgreSQL):} Stores encrypted payload metadata, audit logs, project states, and hybrid encryption keys with arbitrary precision support.
    \item \textbf{On-Chain Settlement Layer (\texttt{GitLockEscrow} Smart Contract on BSC Testnet):} The immutable settlement anchor enforcing fund locking, oracle attestation validation, replay protection, and automated release.
\end{enumerate}

\begin{figure}[H]
\centering
\includegraphics[width=0.92\linewidth]{{"GitLock Secure Code Escrow Architecture"}.png}
\caption{GitLock Secure Code Escrow Architecture and Component Interaction Flow.}
\label{fig:architecture}
\end{figure}

% =========================================================================
% 4. TOOLS AND TECHNOLOGIES USED
% =========================================================================
\section{Tools and Technologies Used}
GitLock integrates industry-standard decentralized and cloud-native engineering frameworks. Table~\ref{tab:tech_stack} summarizes the primary technological components across the stack.

\begin{table}[H]
\centering
\small
\begin{tabularx}{\textwidth}{p{3.2cm} p{3.2cm} X}
\toprule
\textbf{Component / Tool} & \textbf{Classification} & \textbf{Role in GitLock Ecosystem} \\
\midrule
\textbf{Solidity (v0.8.20)} & Smart Contract Language & Implements core escrow logic, ECDSA signature verification, deal state machine, and reentrancy protections. \\
\textbf{Remix IDE} & Web IDE & Used for contract compilation, ABI generation, parameter verification, and Injected Provider testnet deployment. \\
\textbf{Foundry / Hardhat} & Smart Contract Tooling & Local smart contract development, fuzz testing, transaction trace analysis, and gas cost profiling. \\
\textbf{BNB Smart Chain} & Blockchain Network & High-throughput, EVM-compatible Layer-1 network providing low transaction fees and fast settlement times. \\
\textbf{MetaMask} & Web3 Non-Custodial Wallet & Handles client and developer cryptographic identity, transaction authorization, and gas payment on BSC Testnet. \\
\textbf{Next.js 14} & Full-Stack Frontend & Modern React framework utilizing App Router, Tailwind CSS, Lucide icons, and server-side component rendering. \\
\textbf{Express.js / Node.js} & Backend API \& Oracle & Coordinates multi-party notifications, verifies GitHub deployments, manages database state, and signs attestations. \\
\textbf{Ethers.js (v6)} & Web3 JavaScript SDK & Manages on-chain RPC communication, smart contract instantiation, and cryptographic signature computation. \\
\textbf{GitHub REST API} & External Oracular Feed & Programmatically verifies repository existence, commit hash authenticity, deployment statuses, and administrative ownership. \\
\textbf{Google Cloud Run} & Serverless Container Host & Hosts the backend microservice inside a fully managed, auto-scaling stateless container environment with zero downtime. \\
\textbf{Vercel} & Edge Hosting Platform & Hosts the Next.js frontend with global CDN distribution, edge caching, and automated Git-driven deployments. \\
\textbf{Supabase PostgreSQL} & Relational Database & Stores project configurations, arbitrary-precision numerical values (\texttt{NUMERIC}), user roles, and audit trails. \\
\textbf{Google SMTP} & Transactional Email Relay & Dispatches automated notification emails for milestone invitations, deposit confirmations, and transfer alerts. \\
\textbf{GitHub \& Google OAuth} & Authentication Providers & Facilitates frictionless single-sign-on (SSO) and links GitHub developer identities directly to platform accounts. \\
\bottomrule
\end{tabularx}
\caption{Comprehensive Summary of Platform Technologies and Architectural Roles.}
\label{tab:tech_stack}
\end{table}

\subsection{Blockchain Infrastructure Specifications}
The smart contract settlement layer is deployed and verified on the Binance Smart Chain Testnet with the operational specifications outlined in Table~\ref{tab:bsc_specs}.

\begin{table}[H]
\centering
\small
\begin{tabular}{ll}
\toprule
\textbf{Parameter} & \textbf{Value / Configuration} \\
\midrule
Network Name & BNB Smart Chain Testnet (BSC Testnet) \\
Chain ID & \texttt{97} (Hex: \texttt{0x61}) \\
Currency Symbol & \texttt{tBNB} (Testnet BNB) \\
RPC Endpoint & \url{https://data-seed-prebsc-1-s1.binance.org:8545/} \\
Block Explorer & \url{https://testnet.bscscan.com/} \\
Escrow Contract Address & \texttt{0x30CF8a5fa441D39e0604781C24ADeeCBcbfecD49} \\
Backend Oracle Address & \texttt{0x392e227Bf9201A328f9ebD8DBee33f3cBDB15b37} \\
Minimal Test Escrow Amount & \texttt{0.0001 tBNB} \\
Compiler Version & \texttt{Solidity v0.8.20+commit.a1b79de6} \\
Optimization Settings & Enabled (\texttt{200 runs}, EVM Version: \texttt{Paris}) \\
\bottomrule
\end{tabular}
\caption{BNB Smart Chain Testnet Network and Deployment Parameters.}
\label{tab:bsc_specs}
\end{table}

% =========================================================================
% 5. COMPLETE WORKFLOW & PROTOCOL EXECUTION
% =========================================================================
\section{End-to-End Escrow Workflow and the Seven-Step Protocol}
GitLock enforces a deterministic sequence of state transitions across off-chain verification and on-chain settlement, illustrated in Figure~\ref{fig:sequence}.

\begin{figure}[H]
\centering
\includegraphics[width=0.92\linewidth]{{"GitLock Secure Escrow Sequence Diagram"}.png}
\caption{GitLock End-to-End Escrow Protocol Sequence Diagram.}
\label{fig:sequence}
\end{figure}

\subsection{Detailed Protocol Execution (Steps 1 through 7)}
The protocol is executed in seven uninterrupted sequential steps:

\begin{enumerate}
    \item \textbf{Step 1: Project Creation and Developer Milestone Agreement} \\
    The client registers on GitLock, establishes the project parameters, defines the milestone deliverables, and sets the contractual budget in \texttt{tBNB} (e.g., $0.0001\text{ tBNB}$ minimal demo value). An automated invitation is dispatched to the developer via Google SMTP or synced via GitHub OAuth. The Supabase PostgreSQL database persists financial amounts using the unrestricted \texttt{NUMERIC} data type to eliminate fixed-scale decimal truncation.

    \item \textbf{Step 2: Commit Hash Integrity and Staging Deployment Verification} \\
    The developer commits completed code to their GitHub repository. GitLock validates that the commit SHA consists of exactly 40 hexadecimal characters and programmatically verifies via the GitHub REST API that a successful deployment event exists for that exact commit hash:
    \begin{equation}
        \text{Vercel Deployment URL} \equiv \text{Commit SHA Deployment Record}
    \end{equation}
    The system supports both production alias domains (\texttt{app.vercel.app}) and unique hashed preview deployment identifiers.

    \item \textbf{Step 3: Source Code Lock and Hybrid AES-256-CBC Encryption} \\
    The developer exports the repository archive. GitLock computes the SHA-256 tree checksum against the GitHub archive tree to ensure absolute content fidelity. The archive is encrypted using AES-256-CBC with an ephemeral symmetric key, ensuring neither the client nor third parties can view proprietary source code before payment.

    \item \textbf{Step 4: Client Deployment Audit and Smart Contract Escrow Funding} \\
    The client tests the live staging build. Upon client satisfaction, the client approves the deployment and initiates the on-chain funding transaction. Connecting MetaMask to BSC Testnet (Chain ID 97), the client calls:
    \begin{equation}
        \texttt{createAndFundDeal(dealId, developerAddress)} \quad \text{with } \text{msg.value} = 0.0001\text{ tBNB}
    \end{equation}
    The smart contract locks the funds in state \texttt{DealState.Funded}.

    \item \textbf{Step 5: GitHub Repository Ownership Transfer} \\
    Notified via automated email and platform notifications, the developer navigates to the GitHub repository administrative settings and initiates ownership transfer to the client's verified GitHub username (e.g., \texttt{Edavalapati-Aashiq}).

    \item \textbf{Step 6: Off-Chain Oracle Verification and ECDSA Attestation Signing} \\
    The client or developer triggers ownership verification on GitLock. The GitLock Oracle backend queries the GitHub REST API v3 to confirm that the repository's current owner strictly matches the client's GitHub username handle. Once verified, the Oracle signs an EIP-191 compliant ECDSA digest containing the deal ID, developer address, client address, repository name, a 32-byte cryptographically secure random nonce, and a 2-hour deadline.

    \item \textbf{Step 7: On-Chain Payment Release and Source Code Decryption} \\
    The attestation and Oracle signature are submitted to the smart contract via:
    \begin{equation}
        \texttt{releasePaymentWithAttestation(attestation, signature)}
    \end{equation}
    The smart contract recovers the signer address using OpenZeppelin's \texttt{ECDSA.recover}. Confirming the signer is the designated Oracle and the nonce has not been previously consumed, the contract transfers $100\%$ of the locked \texttt{tBNB} directly to the developer's wallet. Simultaneously, GitLock releases the decryption key to the client, unlocking the code archive.
\end{enumerate}

% =========================================================================
% 6. SMART CONTRACT DESIGN
% =========================================================================
\section{Smart Contract Design and Implementation}
The \texttt{GitLockEscrow.sol} contract serves as the decentralized settlement authority. It is implemented in Solidity 0.8.20 and inherits OpenZeppelin's standard \texttt{ReentrancyGuard} to defend against re-entrant calls during native asset transfers.

\subsection{State Machine}
Each escrow deal transitions through four explicit states:
\begin{equation}
    \text{Created} \longrightarrow \text{Funded} \longrightarrow \left\{ \text{Completed} \; \lor \; \text{Refunded} \right\}
\end{equation}

\subsection{Oracle Attestation Verification Logic}
Listing~\ref{lst:escrow_sol} demonstrates the core attestation verification mechanism. The smart contract constructs the Keccak-256 message digest, prepends the standard Ethereum Signed Message prefix, and executes OpenZeppelin's \texttt{ECDSA.recover} to confirm that the signer matches the trusted off-chain Oracle address.

\begin{lstlisting}[language=Solidity, caption={Oracle Attestation Verification in GitLockEscrow.sol}, label={lst:escrow_sol}]
struct Attestation {
    bytes32 dealId;
    address developer;
    address client;
    string repoFullName;
    bytes32 nonce;
    uint256 deadline;
}

function releasePaymentWithAttestation(
    Attestation calldata attestation,
    bytes calldata oracleSignature
) external nonReentrant {
    Deal storage deal = deals[attestation.dealId];
    require(deal.state == DealState.Funded, "Deal is not in Funded state");
    require(block.timestamp <= attestation.deadline, "Attestation expired");
    require(!usedNonces[attestation.nonce], "Nonce already used");
    require(deal.developer == attestation.developer, "Developer mismatch");
    require(deal.client == attestation.client, "Client mismatch");

    // Construct attestation hash
    bytes32 structHash = keccak256(
        abi.encode(
            ATTESTATION_TYPEHASH,
            attestation.dealId,
            attestation.developer,
            attestation.client,
            keccak256(bytes(attestation.repoFullName)),
            attestation.nonce,
            attestation.deadline
        )
    );

    bytes32 ethSignedMessageHash = MessageHashUtils.toEthSignedMessageHash(structHash);
    address recoveredSigner = ECDSA.recover(ethSignedMessageHash, oracleSignature);
    require(recoveredSigner == oracle, "Invalid oracle signature");

    // Prevent replay attacks
    usedNonces[attestation.nonce] = true;
    deal.state = DealState.Completed;

    uint256 paymentAmount = deal.amount;
    (bool success, ) = payable(deal.developer).call{value: paymentAmount}("");
    require(success, "BNB transfer failed");

    emit PaymentReleased(attestation.dealId, deal.developer, paymentAmount);
}
\end{lstlisting}

% =========================================================================
% 7. SECURITY ANALYSIS & THREAT MODEL
% =========================================================================
\section{Security Analysis and Threat Model}
\subsection{Replay Attack Prevention}
Attackers might attempt to intercept a valid Oracle attestation signature from the public blockchain mempool and resubmit it to release funds from a distinct deal. GitLock mitigates this through two complementary mechanisms:
\begin{enumerate}[noitemsep]
    \item \textbf{Deal ID Binding:} The attestation digest tightly couples the transaction to an unalterable \texttt{bytes32 dealId}.
    \item \textbf{Unique Cryptographic Nonce Registry:} The contract maintains a persistent mapping:
    \begin{equation}
        \texttt{mapping(bytes32 => bool) public usedNonces;}
    \end{equation}
    Upon successful execution, \texttt{usedNonces[nonce]} is permanently set to \texttt{true}. Any subsequent invocation with the same nonce reverts immediately.
\end{enumerate}

\subsection{Signature Expiration and Windowing}
To prevent signatures from lingering indefinitely in the event of upstream GitHub transfer anomalies, each attestation specifies a \texttt{deadline} parameter set to two hours ($7,200\text{ seconds}$) from issuance:
\begin{equation}
    \text{Requirement:} \quad \texttt{block.timestamp} \le \texttt{attestation.deadline}
\end{equation}

\subsection{Client Abandonment and Timeout Refunds}
If a client funds the escrow but the developer fails to initiate or complete the repository transfer within a predefined expiration window ($14\text{ days}$), the client can invoke:
\begin{equation}
    \texttt{claimRefundAfterTimeout(bytes32 dealId)}
\end{equation}
This function verifies that $\texttt{block.timestamp} \ge \texttt{deal.fundedAt} + \texttt{TIMEOUT\_PERIOD}$ and transfers $100\%$ of the locked deposit back to the client wallet, preventing permanent asset lockup.

\subsection{Integer Truncation Vulnerability Mitigation}
Early iterations encountered numeric rounding where minimal demo values ($0.0001\text{ tBNB}$) were converted to $0.00$ due to legacy database constraints (\texttt{NUMERIC(10,2)}). The relational schema was modified to arbitrary-precision \texttt{NUMERIC} without scale bounds, ensuring full fidelity down to the Wei level ($10^{-18}\text{ BNB}$).

% =========================================================================
% 8. DEPLOYMENT AND VERIFICATION RESULTS
% =========================================================================
\section{Deployment and Verification Results}
The end-to-end GitLock infrastructure was comprehensively deployed and validated across distributed cloud and blockchain environments.

\subsection{On-Chain Gas Consumption Analysis}
Gas profiling was conducted using Foundry and confirmed via live execution on BSC Testnet. The results are summarized in Table~\ref{tab:gas_analysis}.

\begin{table}[H]
\centering
\small
\begin{tabular}{lrrr}
\toprule
\textbf{Smart Contract Operation} & \textbf{Gas Units (Avg)} & \textbf{Effective Gas Price} & \textbf{Estimated Fee (tBNB)} \\
\midrule
Contract Deployment & $1,248,510$ & $3.0\text{ Gwei}$ & $0.003745\text{ tBNB}$ \\
\texttt{createAndFundDeal()} & $88,432$ & $3.0\text{ Gwei}$ & $0.000265\text{ tBNB}$ \\
\texttt{releasePaymentWithAttestation()} & $64,219$ & $3.0\text{ Gwei}$ & $0.000192\text{ tBNB}$ \\
\texttt{claimRefundAfterTimeout()} & $41,120$ & $3.0\text{ Gwei}$ & $0.000123\text{ tBNB}$ \\
\bottomrule
\end{tabular}
\caption{Smart Contract Execution Gas Benchmarks on BSC Testnet.}
\label{tab:gas_analysis}
\end{table}

\subsection{Live Verification Proof}
The complete live cycle was verified with actual production parameters:
\begin{itemize}[noitemsep]
    \item \textbf{Repository Under Test:} \texttt{Aashiq-Edavalapati/gitlock-demo}
    \item \textbf{Commit Hash:} \texttt{8e81cfa6622276d919d1be0db4d32c19664a6ee6}
    \item \textbf{Verified Staging URL:} \url{https://gitlock-demo.vercel.app/}
    \item \textbf{Transferred Owner (Client):} \texttt{Edavalapati-Aashiq}
    \item \textbf{Escrow Deal ID:} \texttt{0x689af73388b3760388bf9f7a539d291783a2659acdbbbc7c610019a1bdb5a189}
    \item \textbf{Escrow Amount:} $0.0001\text{ tBNB}$
    \item \textbf{Oracle Attestation Status:} Validated, signed, and broadcasted to BSC Testnet.
\end{itemize}

% =========================================================================
% 9. FUTURE IMPROVEMENTS
% =========================================================================
\section{Future Improvements}
While GitLock successfully provides trustless escrow for software engineering, future development trajectories will expand platform decentralization:
\begin{enumerate}
    \item \textbf{Decentralized Oracle Networks (Chainlink Functions):} Transitioning from a single backend Oracle signer to a decentralized consensus network using Chainlink Functions or UMA's Optimistic Oracle to remove single-operator trust assumptions.
    \item \textbf{Zero-Knowledge Proof of Compilation (zk-SNARKs):} Implementing RISC Zero or Succinct SP1 zkVMs to generate cryptographic proofs that a private source code archive compiles cleanly and satisfies a suite of automated unit tests without revealing the proprietary source code before payment.
    \item \textbf{Multi-Chain Interoperability:} Deploying LayerZero or Chainlink CCIP cross-chain bridges to enable clients to fund escrow using assets on Ethereum, Arbitrum, or Polygon while settling payments to developers on low-cost networks.
    \item \textbf{Decentralized Autonomous Arbitration (DAO / Kleros Integration):} Providing subjective dispute resolution for non-binary milestone disagreements (such as UI design aesthetics or edge-case software bugs) through crowdsourced economic incentives via Kleros or Aragon Court.
\end{enumerate}

% =========================================================================
% 10. CONCLUSION
% =========================================================================
\section{Conclusion}
This case study presented the design, implementation, and empirical verification of \textbf{GitLock}, a blockchain-anchored, trustless code delivery and escrow protocol. By synthesizing non-custodial smart contracts on the BNB Smart Chain with off-chain cryptographic Oracle verification of GitHub repository ownership, GitLock successfully resolves the fundamental counterparty dilemma in remote software engineering. The platform eliminates centralized commission fees, guarantees financial atomicity, prevents code theft through hybrid client-side encryption, and provides transparent on-chain auditability. Live deployments across Vercel, Google Cloud Run, Supabase, and the BSC Testnet confirm the practical viability, robust security, and economic efficiency of programmatic escrow protocols for the future of digital software commerce.

% =========================================================================
% REFERENCES
% =========================================================================
\begin{thebibliography}{99}
\bibitem{nakamoto2008}
S.~Nakamoto, ``Bitcoin: A Peer-to-Peer Electronic Cash System,'' \textit{Decentralized Business Review}, p.~21260, 2008.

\bibitem{buterin2014}
V.~Buterin, ``A Next-Generation Smart Contract and Decentralized Application Platform,'' \textit{Ethereum White Paper}, vol.~3, no.~37, pp.~2--1, 2014.

\bibitem{openzeppelin}
OpenZeppelin, ``OpenZeppelin Contracts: Secure Smart Contract Library in Solidity,'' 2024. [Online]. Available: \url{https://docs.openzeppelin.com/contracts/}

\bibitem{eip712}
M.~Cerrador, W.~Schwab, et~al., ``EIP-712: Typed Structured Data Hashing and Signing,'' \textit{Ethereum Improvement Proposals}, no.~712, 2018.

\bibitem{bsc2020}
BNB Chain Development Team, ``BNB Smart Chain (BSC) Whitepaper: A Dual-Chain Architecture for Programmability and High Performance,'' 2020. [Online]. Available: \url{https://docs.bnbchain.org/}

\bibitem{ethersjs}
R.~Moore, ``Ethers.js: Complete, Compact, and Simple Library for Ethereum and EVM-Compatible Blockchains,'' v6 Documentation, 2024. [Online]. Available: \url{https://docs.ethers.org/v6/}

\bibitem{githubapi}
GitHub Inc., ``GitHub REST API Documentation: Repositories and Deployments Reference,'' 2024. [Online]. Available: \url{https://docs.github.com/en/rest}
\end{thebibliography}

\end{document}
